% Source: https://artofproblemsolving.com/community/c3h2045167_the_smog_of_the_log
% Stable identity: aops-topic-2045167-q01
\textbf{(a)} Take \(a=m\) and \(b=n\). Then \(\log_m a+\log_n b=1+1=2\), an integer. (More generally \(a=m^p\), \(b=n^q\) with nonnegative integers \(p,q\) gives the integer \(p+q\).) So the answer to the general question is \emph{yes}, sometimes.

\textbf{(b)} Since \(6=2\cdot 3\),
\[
\log_2 6+\log_3 6=(1+\log_2 3)+(1+\log_3 2)=2+t+\dfrac{1}{t},\qquad t=\log_2 3 .
\]
Because \(2<3<4\), we have \(1<t<2\). The function \(g(t)=t+\dfrac{1}{t}\) is increasing for \(t>1\) (indeed \(g(t)-g(s)=(t-s)\left(1-\dfrac{1}{ts}\right)>0\) for \(t>s\ge1\)), so
\[
2=g(1)<t+\dfrac{1}{t}<g(2)=\dfrac{5}{2} .
\]
Hence \(4<\log_2 6+\log_3 6<\dfrac{9}{2}\). There is no integer strictly between \(4\) and \(4.5\), so \(\log_2 6+\log_3 6\) is not an integer. \(\blacksquare\)
